\documentclass[12pt, a4paper]{article}

% ── PACKAGES ──────────────────────────────────────────────────
\usepackage[utf8]{inputenc}
\usepackage[T1]{fontenc}
\usepackage{geometry}
\geometry{margin=2.5cm}
\usepackage{hyperref}
\hypersetup{
    colorlinks=true,
    linkcolor=blue,
    urlcolor=blue,
    citecolor=blue,
    pdftitle={CDKN1A (p21) Level as a Quantitative
              Predictor of CDK4/6 Inhibitor Benefit
              Magnitude in Luminal A Breast Cancer},
    pdfauthor={Eric Robert Lawson},
    pdfkeywords={CDKN1A, p21, CDK4/6 inhibitor,
                 palbociclib, ribociclib, abemaciclib,
                 luminal A, breast cancer, biomarker,
                 attractor geometry, Waddington landscape,
                 MONARCH, PALOMA, MONALEESA,
                 OrganismCore, cell cycle, endocrine therapy}
}
\usepackage{amsmath}
\usepackage{booktabs}
\usepackage{longtable}
\usepackage{array}
\usepackage{parskip}
\usepackage{microtype}
\usepackage{authblk}
\usepackage{titlesec}
\usepackage{fancyhdr}
\usepackage{xcolor}
\usepackage{float}
\usepackage{enumitem}
\usepackage{setspace}

% ── COLOURS ───────────────────────────────────────────────────
\definecolor{repobox}{RGB}{240,247,255}
\definecolor{findingbox}{RGB}{245,255,245}
\definecolor{novelbox}{RGB}{255,248,220}
\definecolor{mechanismbox}{RGB}{250,245,255}
\definecolor{actionbox}{RGB}{230,255,230}
\definecolor{warnbox}{RGB}{255,235,235}

% ── HEADER / FOOTER ───────────────────────────────────────────
\pagestyle{fancy}
\fancyhf{}
\rhead{\small OrganismCore \textbar{} 2026-03-06}
\lhead{\small CS-LIT-14: CDKN1A as Quantitative CDK4/6i
       Benefit Predictor in LumA}
\rfoot{\small Page \thepage}
\renewcommand{\headrulewidth}{0.4pt}

% ── SECTION FORMATTING ────────────────────────────────────────
\titleformat{\section}
  {\large\bfseries}{\thesection}{1em}{}[\titlerule]
\titleformat{\subsection}
  {\normalsize\bfseries}{\thesubsection}{1em}{}
\titleformat{\subsubsection}
  {\small\bfseries}{\thesubsubsection}{1em}{}

% ══════════════════════════════════════════════════════════════
% TITLE
% ══════════════════════════════════════════════════════════════
\title{
    \vspace{-1em}
    \textbf{CDKN1A (p21) Level as a Quantitative
    Predictor of CDK4/6 Inhibitor Benefit Magnitude
    in Luminal A Breast Cancer:}\\[0.5em]
    \large A Geometry-Derived Biomarker Hypothesis
    with Mechanistic Grounding and Proposed
    Validation in MONARCH, PALOMA, and
    MONALEESA Trial Datasets\\[0.4em]
    \normalsize \textit{OrganismCore Framework ---
    Document CS-LIT-14}\\[0.3em]
    \small Derived from Waddington Attractor Geometry
    Analysis of 19,542 Single Cancer Cells (GSE176078)\\[0.3em]
    \small Part of the Six Lock Type Classification
    of Breast Cancer (CS-LIT-2, Zenodo 2026-03-06)
}

\author[1]{Eric Robert Lawson}
\affil[1]{
    OrganismCore\\
    \href{mailto:OrganismCore@proton.me}
         {OrganismCore@proton.me}\\
    ORCID:
    \href{https://orcid.org/0009-0002-0414-6544}
         {0009-0002-0414-6544}
}

\date{March 6, 2026}

% ══════════════════════════════════════════════════════════════
\begin{document}
% ══════════════════════════════════════════════════════════════

\maketitle
\thispagestyle{fancy}

% ── REPOSITORY POINTER ────────────────────────────────────────
\begin{center}
\colorbox{repobox}{
\begin{minipage}{0.88\textwidth}
\vspace{0.5em}
\centering
\textbf{Full computational record --- scripts, data caches,
reasoning artifacts, and raw logs:}\\[0.4em]
\href{https://github.com/Eric-Robert-Lawson/attractor-oncology}
{\texttt{https://github.com/Eric-Robert-Lawson/attractor-oncology}}
\\[0.3em]
\small Part of the OrganismCore BRCA Cross-Subtype
Analysis Series (BRCA-S8h, 2026-03-05).\\
All predictions were locked from attractor geometry
\textbf{before} any literature review was performed.
\vspace{0.5em}
\end{minipage}}
\end{center}

\vspace{1em}

% ── NOVEL NOTICE ──────────────────────────────────────────────
\begin{center}
\colorbox{novelbox}{
\begin{minipage}{0.88\textwidth}
\vspace{0.5em}
\textbf{VERDICT: NOVEL}\\[0.3em]
This finding has no direct published equivalent as of
2026-03-06. CDKN1A loss is established as a resistance
mechanism; the specific hypothesis advanced here---that
\emph{baseline quantitative p21 level predicts the
magnitude of benefit}, not merely binary sensitivity or
resistance---has not been formally published or
prospectively validated. This document establishes
priority and proposes the validation pathway.
\vspace{0.5em}
\end{minipage}}
\end{center}

\vspace{1em}

% ════════════════════════════════════���═════════════════════════
\begin{abstract}
% ══════════════════════════════════════════════════════════════
Luminal A breast cancer is the most prevalent breast cancer
subtype and the primary population for CDK4/6 inhibitor
therapy. Current clinical use of palbociclib, ribociclib, and
abemaciclib does not stratify patients by the magnitude of
expected benefit --- all HR+/HER2$-$ patients are treated as
a uniform population. We report a geometry-derived hypothesis,
originating from Waddington attractor analysis of 19,542
single cancer cells (GSE176078), that baseline quantitative
CDKN1A (p21) expression level predicts not merely whether a
Luminal A patient will respond to CDK4/6 inhibition, but
\emph{how much benefit they will derive}. The geometric
basis is as follows: in Luminal A, CDKN1A loss is the
primary identified lock mechanism --- the attractor is
maintained by CDK-mediated suppression of differentiation
programs, and CDKN1A is the quantitative gate on that
suppression. Higher residual p21 means the lock is
shallower and more susceptible to CDK4/6 inhibition;
lower p21 means the lock is deeper and the CDK4/6i
effect is correspondingly attenuated. This produces a
continuous dose--response relationship between p21 level
and CDK4/6i benefit magnitude that is not captured by
binary resistance/sensitivity classification. We propose
this hypothesis as a specific, falsifiable, quantitative
prediction testable in existing trial biobank datasets
from MONARCH, PALOMA, and MONALEESA, requiring no new
patient recruitment.
\end{abstract}

\tableofcontents
\newpage

% ══════════════════════════════════════════════════════════════
\section{Context and Origin}
% ══════════════════════════════════════════════════════���═══════

This document is part of the OrganismCore BRCA Cross-Subtype
Analysis (Document BRCA-S8h, 2026-03-05), which applied
Waddington attractor geometry to 19,542 single cancer cells
across all six major breast cancer subtypes (GSE176078).

The analysis produced 30 pre-specified items subjected to
a post-hoc literature check. CS-LIT-14 was recorded as a
NOVEL finding --- no published equivalent for the quantitative
continuous-predictor formulation was identified.

The parent document establishing the six-lock taxonomy within
which this finding sits is:

\begin{itemize}
    \item \textbf{CS-LIT-2}: A Six Lock Type Classification
    of Breast Cancer (Zenodo, 2026-03-06,
    \href{https://doi.org/10.5281/zenodo.18884157}
    {DOI: 10.5281/zenodo.18884157})
\end{itemize}

All findings in this series were derived from attractor
geometry \textbf{before} any literature review was
performed. The prediction was locked. The literature was
checked afterward. The NOVEL verdict means no published
equivalent was found at the time of the check.

% ══════════════════════════════════════════════════════════════
\section{The Luminal A Lock --- Geometric Basis}
% ══════════════════════════════════════════════════════════════

\subsection{What Attractor Geometry Reveals in Luminal A}

Luminal A occupies the shallowest attractor position in the
FOXA1/EZH2 ordering axis across all six breast cancer subtypes
(FOXA1/EZH2 ratio = 9.38, the highest of any subtype).
This position reflects high FOXA1 expression and low EZH2
activity --- the identity programs of the luminal lineage are
substantially intact. The Waddington valley is shallow and
relatively accessible.

The question the geometry raises is: if the attractor is
already shallow, what is holding the cell in a cancerous
state? What is the specific lock mechanism in Luminal A?

The answer derived from the attractor geometry is
\textbf{CDK-mediated cell cycle dysregulation as the primary
maintaining force}. In Luminal A, the chromatin programs
of luminal identity are largely present. What is disrupted
is the cell cycle gate --- specifically, the suppression of
CDKN1A (p21), which functions as the quantitative brake on
CDK4/6-mediated proliferative drive.

\subsection{CDKN1A as the Quantitative Gate}

In normal luminal epithelium, p21 imposes context-dependent
cell cycle arrest, allowing differentiation programs to
consolidate. In Luminal A cancer, p21 is suppressed or
functionally attenuated --- not necessarily through mutation
or complete loss, but through quantitative downregulation
sufficient to allow persistent CDK4/6 activity to drive
proliferation.

This creates a \textbf{continuous relationship} between
p21 level and lock depth:

\begin{center}
\colorbox{mechanismbox}{
\begin{minipage}{0.82\textwidth}
\vspace{0.5em}
\textbf{Geometric mechanism:}\\[0.4em]
High p21 (CDKN1A) $\rightarrow$ shallow lock $\rightarrow$
CDK4/6i strongly effective \\[0.3em]
Intermediate p21 $\rightarrow$ intermediate lock $\rightarrow$
CDK4/6i moderately effective \\[0.3em]
Low p21 $\rightarrow$ deep lock via alternative CDK pathway
engagement $\rightarrow$ CDK4/6i attenuated \\[0.3em]
Very low / absent p21 $\rightarrow$ CDK2 bypass available
$\rightarrow$ CDK4/6i resistance
\vspace{0.5em}
\end{minipage}}
\end{center}

\vspace{0.5em}

The established literature treats CDKN1A loss as a
\textbf{binary resistance mechanism} --- present or absent.
This framework predicts a \textbf{quantitative continuous
relationship}: baseline p21 level predicts the
\emph{magnitude} of CDK4/6i benefit across the full range
of expression, not merely the binary resistant/sensitive
boundary.

% ══════════════════════════════════════════════════════════════
\section{The Precise Hypothesis}
% ══════════════════════════════════════════════════════════════

\begin{center}
\colorbox{findingbox}{
\begin{minipage}{0.88\textwidth}
\vspace{0.6em}
\textbf{The specific, falsifiable prediction:}\\[0.5em]
In Luminal A HR+/HER2$-$ breast cancer patients treated
with a CDK4/6 inhibitor plus endocrine therapy, baseline
quantitative CDKN1A (p21) expression level --- measured
as a continuous variable, not a binary cutoff --- will
correlate positively and monotonically with progression-free
survival benefit from CDK4/6 inhibition.\\[0.4em]
Specifically:
\begin{enumerate}[leftmargin=*, itemsep=3pt]
    \item Patients in the highest quartile of baseline
    CDKN1A expression will derive the largest PFS benefit
    from CDK4/6i addition to endocrine therapy.
    \item Patients in the lowest quartile of baseline
    CDKN1A expression will derive the smallest PFS benefit,
    approaching the endocrine-therapy-alone arm.
    \item The relationship will be monotonic across
    quartiles, not threshold-dependent.
    \item This gradient will be most pronounced in
    Luminal A (PAM50 LumA) compared with Luminal B
    (PAM50 LumB), because the CDK lock is the
    \emph{primary} mechanism in LumA and a
    \emph{secondary} mechanism in LumB.
\end{enumerate}
\vspace{0.4em}
\end{minipage}}
\end{center}

% ══════════════════════════════════════════════════════════════
\section{What Distinguishes This From the Existing Literature}
% ══════════════════════════════════════════════════════════════

\subsection{What the Literature Currently Contains}

The existing literature on CDKN1A and CDK4/6 inhibitors
addresses three distinct claims:

\begin{enumerate}
    \item \textbf{CDKN1A loss causes resistance.}
    This is established. Complete or near-complete loss
    of p21 allows CDK2 to compensate for CDK4/6
    inhibition, maintaining cell cycle progression.
    This is a binary resistance mechanism.

    \item \textbf{p21 upregulation is a pharmacodynamic
    marker of CDK4/6i activity.}
    This is established. Entinostat and CDK4/6 inhibitors
    both upregulate p21 as a downstream effect.
    Post-treatment p21 level reflects drug action.

    \item \textbf{p21 restoration re-sensitises resistant
    cells.}
    This is established in preclinical models.
    Restoring p21 expression recovers CDK4/6i
    sensitivity.
\end{enumerate}

\subsection{What the Literature Does Not Contain}

None of the above three claims is equivalent to the
hypothesis advanced here. This document proposes:

\begin{itemize}
    \item \textbf{Baseline} p21 level (pre-treatment,
    not post-treatment pharmacodynamic response)
    \item As a \textbf{continuous quantitative predictor}
    (not binary resistance/sensitivity)
    \item Of \textbf{benefit magnitude} (not just
    whether a response occurs)
    \item Specifically in \textbf{Luminal A} as a
    subtype-stratified prediction
    \item Testable as a \textbf{monotonic gradient across
    expression quartiles}
\end{itemize}

This is a different and more specific claim. It is
falsifiable in existing trial biobank data without new
patient recruitment.

% ══════════════════════════════════════════════════════════════
\section{Mechanistic Grounding}
% ═════════════════════════════════════════════════════════���════

\subsection{Why the Relationship Should Be Continuous,
Not Binary}

CDK4/6 inhibitors work by releasing p21 and p27 from
sequestration by cyclin D--CDK4/6 complexes, allowing
these inhibitors to block Rb phosphorylation and arrest
the cell cycle at G1. The efficiency of this arrest is
directly proportional to the available pool of functional
p21. A cell with high baseline p21 has more arrest
capacity available to be engaged by CDK4/6 inhibition.
A cell with low but non-zero p21 has less. A cell with
near-zero p21 has minimal arrest capacity, and CDK2
can substitute.

This is not a switch. It is a gradient. The geometry
predicts the clinical consequence of that gradient:
a monotonic relationship between baseline p21 and
CDK4/6i benefit.

\subsection{Why This Is Subtype-Specific to Luminal A}

In Luminal B, the attractor is maintained by multiple
co-active mechanisms --- the DNMT3A/HDAC2 chromatin lock
on ER output (CS-LIT-8/9) operates in parallel with CDK
dysregulation. The CDK axis is not the sole or primary
lock. Consequently, p21 level in LumB is a less clean
predictor of CDK4/6i benefit because other lock
mechanisms can maintain the attractor even when the
CDK gate is partially restored.

In Luminal A, the CDK axis \emph{is} the primary lock.
The chromatin programs are intact. The cell cycle gate is
the specific vulnerability. p21 level in LumA therefore
maps more directly and cleanly onto CDK4/6i benefit
than in any other subtype.

This subtype specificity is a testable prediction:
the CDKN1A--benefit gradient should be steeper in
PAM50 LumA than in PAM50 LumB when both populations
are analysed in the same trial dataset.

% ══════════════════════════════════════════════════════════════
\section{Proposed Validation Pathway}
% ══════════════════════════════════════════════════════════════

\subsection{Why No New Patients Are Required}

The MONARCH, PALOMA, and MONALEESA trial series enrolled
thousands of HR+/HER2$-$ patients with paired baseline
tumor biopsies and detailed outcome data. Correlative
biomarker analyses are standard in these trials. The
biobank material to test this hypothesis exists.

\subsection{Specific Validation Design}

\begin{center}
\colorbox{actionbox}{
\begin{minipage}{0.88\textwidth}
\vspace{0.5em}
\textbf{Proposed analysis (requires no new recruitment):}
\begin{enumerate}[leftmargin=*, itemsep=4pt]
    \item \textbf{Dataset:} Baseline tumor RNA-seq or
    qRT-PCR from any MONARCH, PALOMA, or MONALEESA
    trial biobank with available PAM50 subtyping and
    PFS outcome data.
    \item \textbf{Population:} Restrict to PAM50 Luminal
    A patients receiving CDK4/6i + endocrine therapy
    versus endocrine therapy alone.
    \item \textbf{Measurement:} Baseline CDKN1A
    expression as a continuous variable (normalised
    mRNA counts or IHC H-score).
    \item \textbf{Analysis:} Divide CDKN1A into
    quartiles. Compare PFS benefit of CDK4/6i addition
    across quartiles using interaction term in Cox
    proportional hazards model (CDK4/6i arm $\times$
    CDKN1A quartile).
    \item \textbf{Primary endpoint:} Monotonic gradient
    in CDK4/6i benefit across CDKN1A quartiles
    (Q1 $<$ Q2 $<$ Q3 $<$ Q4).
    \item \textbf{Secondary endpoint:} Steeper gradient
    in PAM50 LumA versus PAM50 LumB (subtype
    interaction test).
    \item \textbf{Falsification criterion:} A flat or
    non-monotonic relationship across quartiles falsifies
    the quantitative continuous prediction. A U-shaped
    relationship would falsify the directional hypothesis.
\end{enumerate}
\vspace{0.4em}
\end{minipage}}
\end{center}

\subsection{What a Positive Result Would Mean Clinically}

If the monotonic gradient is confirmed:

\begin{itemize}
    \item Baseline CDKN1A measurement (by IHC or RNA)
    becomes a patient stratification tool for CDK4/6i
    prescribing in Luminal A.
    \item Patients with high p21 would be identified as
    the subgroup deriving the largest absolute benefit,
    strengthening the case for CDK4/6i in that group.
    \item Patients with very low p21 would be identified
    as the subgroup deriving minimal benefit from CDK4/6i,
    redirecting clinical attention toward alternative
    lock-breaking strategies for that population.
    \item The measurement is a standard RNA expression
    value or a simple IHC stain --- implementable in any
    clinical laboratory.
\end{itemize}

\subsection{Relevant Trial Contacts}

The hypothesis is most directly testable in datasets
controlled by the following groups:

\begin{itemize}
    \item \textbf{MONARCH series (abemaciclib):}
    Eli Lilly oncology biomarker research group;
    principal investigators of MONARCH-2 and MONARCH-3.
    \item \textbf{PALOMA series (palbociclib):}
    Pfizer oncology; principal investigators of
    PALOMA-2 and PALOMA-3 correlative science arms.
    \item \textbf{MONALEESA series (ribociclib):}
    Novartis oncology; principal investigators of
    MONALEESA-2, MONALEESA-3, and MONALEESA-7.
\end{itemize}

% ══════════════════════════════════════════════════════════════
\section{Relationship to the Six Lock Type Classification}
% ══════════════════════════════════════════════════════════════

This finding is the drug prediction derived from
Lock Type 1 (Luminal A) in the six-lock taxonomy
(CS-LIT-2). The complete taxonomy is:

\begin{center}
\begin{tabular}{lll}
\toprule
\textbf{Subtype} & \textbf{Lock Type} &
\textbf{Primary Drug Target} \\
\midrule
Luminal A & CDK-mediated gate (CDKN1A loss) &
CDK4/6 inhibitors \\
Luminal B & DNMT3A/HDAC2 chromatin lock &
HDAC inhibitors (entinostat) \\
HER2-enriched & EZH2 lock (deep fraction) &
EZH2i + anti-HER2 \\
TNBC & EZH2 PRC2 epigenetic lock &
Tazemetostat \\
Claudin-low & Structural immune lock &
Anti-TIGIT sequence \\
ILC & Structural FOXA1 lock (inverse) &
Fulvestrant over AI \\
\bottomrule
\end{tabular}
\end{center}

\vspace{0.5em}

CS-LIT-14 is the quantitative refinement of the
Luminal A drug prediction: not merely that CDK4/6
inhibitors are the correct drug class, but that
\emph{baseline p21 level predicts the magnitude of
benefit within that drug class}.

% ═════════════════════════════════════��══════════���═════════════
\section{What This Document Claims and Does Not Claim}
% ══════════════════════════════════════════════════════════════

\subsection{What IS claimed}

\begin{itemize}
    \item The geometric basis for a continuous
    quantitative relationship between baseline CDKN1A
    level and CDK4/6i benefit magnitude in Luminal A
    is mechanistically grounded and geometrically derived.
    \item This specific quantitative formulation has not
    been published as a formal hypothesis with a
    prospective validation design.
    \item The hypothesis is falsifiable with existing
    trial biobank data requiring no new patient
    recruitment.
    \item The subtype specificity to Luminal A
    (versus Luminal B) is a testable and distinguishing
    prediction.
\end{itemize}

\subsection{What is NOT claimed}

\begin{itemize}
    \item This has not been validated in any dataset.
    No clinical data are presented in this document.
    \item The quantitative cutoffs for clinically
    meaningful p21 strata have not been determined.
    \item This does not replace or supersede the
    established binary resistance literature. It extends
    it into a continuous prediction that the binary
    literature does not address.
    \item Clinical implementation requires prospective
    validation.
\end{itemize}

% ══════════════════════════════════════════════════════════════
\section{Locked Statement}
% ══════════════════════════════════════════════════════════════

\begin{center}
\colorbox{findingbox}{
\begin{minipage}{0.88\textwidth}
\vspace{0.6em}
\textbf{Stated once, precisely, without inflation:}\\[0.5em]
Waddington attractor geometry applied to 19,542 single
breast cancer cells (GSE176078) identifies CDKN1A loss
as the primary lock mechanism in Luminal A breast cancer.
This geometry predicts that baseline quantitative p21
level will function as a continuous predictor of CDK4/6
inhibitor benefit magnitude --- not merely as a binary
resistance marker --- with the largest benefit in
p21-high patients and the smallest in p21-low patients.
This prediction is specific to Luminal A, mechanistically
grounded, falsifiable in existing trial datasets, and
has no direct published equivalent as of 2026-03-06.
\vspace{0.5em}
\end{minipage}}
\end{center}

% ══════════════════════════════════════════════════════════════
\section{Document Metadata}
% ══════════════════════════════════════════════════════════════

\begin{tabular}{ll}
\toprule
\textbf{Field} & \textbf{Value} \\
\midrule
Document ID & CS-LIT-14 \\
Parent document & BRCA-S8h (2026-03-05) \\
Verdict & NOVEL \\
Date & 2026-03-06 \\
Author & Eric Robert Lawson \\
ORCID & 0009-0002-0414-6544 \\
Contact & OrganismCore@proton.me \\
Repository &
\href{https://github.com/Eric-Robert-Lawson/attractor-oncology}
{attractor-oncology} \\
Source data & GSE176078 (19,542 single cancer cells) \\
Parent DOI &
\href{https://doi.org/10.5281/zenodo.18884157}
{10.5281/zenodo.18884157} (CS-LIT-2) \\
License & CC BY 4.0 \\
\bottomrule
\end{tabular}

\end{document}
